% =====================================================================
%  CAPÍTULO 1 · MARCO TEÓRICO
% =====================================================================
\chapter{Marco teórico}
\label{cap:marco}

Este capítulo reúne los conceptos, enfoques y antecedentes necesarios para
comprender el tema. Reemplace el texto de ejemplo por su propio contenido.

\section{Antecedentes}
\label{sec:antecedentes}

Resuma los estudios previos más relevantes y señale qué vacío o discusión
dejan abierta. Cuando cite textualmente, use comillas y el número de página:
\citep[p.~25]{eco2001tesis}.

\section{Conceptos fundamentales}
\label{sec:conceptos}

\subsection{Primer concepto}

Defina el concepto con precisión y acompáñelo de su fuente
\citep{hernandez2014metodologia}.

\subsection{Segundo concepto}

Explique el segundo concepto y su relación con el primero. Para enumerar
elementos puede usar listas:

\begin{itemize}
  \item Elemento uno, con una breve explicación.
  \item Elemento dos, con una breve explicación.
  \item Elemento tres, con una breve explicación.
\end{itemize}

\section{Enfoques teóricos}
\label{sec:enfoques}

Compare las principales posturas. La Tabla~\ref{tab:enfoques} ilustra cómo
presentar una comparación de manera ordenada.

\begin{table}[htbp]
  \centering
  \caption{Comparación de enfoques teóricos (ejemplo).}
  \label{tab:enfoques}
  \begin{tabular}{lll}
    \toprule
    Enfoque & Idea central & Autores representativos \\
    \midrule
    Enfoque A & Descripción breve & Autor 1 (año) \\
    Enfoque B & Descripción breve & Autor 2 (año) \\
    Enfoque C & Descripción breve & Autor 3 (año) \\
    \bottomrule
  \end{tabular}
  \par\smallskip\footnotesize Fuente: elaboración propia.
\end{table}
